\section{Two actual oriented profiles in a quadratic-field cover}
\label{dc-section}

The actual quartic norm admits a more direct lifting than the generic
splitting-field construction of Sections~\ref{qc-section}
and~\ref{mx-section}. Its defining integer is itself a primitive
Eisenstein norm. This lets us use two actual residual coefficients,
with no coefficient-height term proportional to either exponent.
The resulting point-height problem remains open.

Retain $\mathcal O=\mathbb Z[\zeta]$, $\zeta^2-\zeta+1=0$,
$K_0=\mathbb Q(\zeta)$ and $\gamma=1+\zeta$. For positive coprime
integers $a,b$, put
\[
 c=a+b,\quad x=a/b,\quad M=a^2+ab+b^2,
 \quad F=a^4+3a^3b+5a^2b^2+3ab^3+b^4.
\]
Assume an actual first oriented profile
\begin{equation}\label{dc-first-profile}
 a+b\zeta=u_0\gamma^e v_0w_0^h,\qquad
 M=3^eV_0R^h,\qquad
 V_0=\operatorname N v_0,\quad R=\operatorname N w_0\ge7,
\end{equation}
where $u_0$ is a unit, $e\in\{0,1\}$, $h\ge2$, and
$v_0,w_0$ are primitive oriented unramified elements, as in
\eqref{lr-profile}. Their oriented supports may overlap.
For the second norm assume only
\begin{equation}\label{dc-second-profile}
 F=V_1Q^g,\qquad g\ge2,\qquad V_1,Q\in\mathbb Z_{>0},\qquad Q>1.
\end{equation}
All exponents are integers. The residual $V_1$ need not be $g$-free;
that is an optional canonical convention. The exponents need not
have a common divisor, or be coprime to six. As checked in
Section~\ref{mx-section}, $2,3,5$ do not divide $F$, so $Q\ge7$.

\begin{lemma}[An actual second oriented factorization]\label{dc-second-lift}
There are a unit $u_1$ and primitive oriented unramified elements
$v_1,w_1\in\mathcal O$ such that
\begin{equation}\label{dc-actual-second}
 ab+M\zeta=u_1v_1w_1^g,\qquad
 \operatorname N v_1=V_1,\qquad\operatorname N w_1=Q.
\end{equation}
\end{lemma}
\begin{proof}
Direct expansion gives
\begin{equation}\label{dc-second-norm-identity}
 \operatorname N(ab+M\zeta)=(ab)^2+abM+M^2=F.
\end{equation}
Moreover $\gcd(ab,M)=1$: a prime dividing $a$ or $b$ and $M$
would divide both $a$ and $b$. Thus $ab+M\zeta$ is primitive,
and it is unramified because $3\nmid F$.

Use unique factorization in the Euclidean Eisenstein ring. An inert
rational prime in the norm of a primitive element would divide both
coordinates, as would the simultaneous occurrence of conjugate split
factors. There is therefore a factorization
\[
 ab+M\zeta=u_1\prod_q\pi_q^{E_q},\qquad
 \operatorname N\pi_q=q,
\]
with precisely one orientation chosen at each rational prime in the
support. Equation~\eqref{dc-second-profile} gives
$E_q=v_q(V_1)+g v_q(Q)$, with both summands nonnegative. Set
\[
 v_1=\prod_q\pi_q^{v_q(V_1)},\qquad
 w_1=\prod_q\pi_q^{v_q(Q)}.
\]
Their product and norms give \eqref{dc-actual-second} exactly.
They inherit primitivity, unramifiedness and the absence of
conjugate pairs within either factor. They may share oriented
factors. No bound $v_q(V_1)<g$ is used.
\end{proof}

This is an existence proof using unique factorization, not an assertion
that a practical factorization of a large $F$ has already been computed.

Put $N(X)=X^2+X+1$, $\beta_+=-\zeta$, $\beta_-=-\bar\zeta$,
and
\begin{equation}\label{dc-quadratics}
 G(X)=X+\zeta N(X),\qquad
 \bar G(X)=X+\bar\zeta N(X).
\end{equation}
Define, in the fixed quadratic field $K_0$,
\begin{equation}\label{dc-coefficients}
 \eta_0=(u_0^2\zeta^e v_0/\bar v_0)^{-1},\qquad
 \eta_1=(u_1^2v_1/\bar v_1)^{-1}.
\end{equation}

\begin{theorem}[A degree-$hg$ cover from two actual profiles]
\label{dc-cover}
The actual profiles yield nonzero $K_0$-valued coordinates satisfying
\begin{equation}\label{dc-kummer-equations}
 Y^h=\eta_0\frac{x-\beta_+}{x-\beta_-},\qquad
 Z^g=\eta_1\frac{G(x)}{\bar G(x)},\qquad
 Y=w_0/\bar w_0,\quad Z=w_1/\bar w_1.
\end{equation}
Their coefficients have exact absolute logarithmic Weil heights
\begin{equation}\label{dc-exact-height}
 h_{K_0}(\eta_0)=\tfrac12\log V_0,\qquad
 h_{K_0}(\eta_1)=\tfrac12\log V_1.
\end{equation}
The smooth projective normalization is geometrically connected over
$K_0$, with degree $D=hg$ over $\mathbb P^1_x$ and genus
\begin{equation}\label{dc-genus}
 G_C=1+hg\left(2-\frac1h-\frac2g\right)
     =1+2hg-g-2h.
\end{equation}
Every actual profile gives a $K_0$-point of this curve above $a/b$.
\end{theorem}
\begin{proof}
The first equation follows by conjugating the first profile, using
$\gamma/\bar\gamma=\zeta$ and $u_0/\bar u_0=u_0^2$.
For the second, $(ab+M\zeta)/b^2=G(x)$; divide
\eqref{dc-actual-second} by its conjugate and rearrange.
Neither step takes an exponent-th root of a unit.
Lemma~\ref{lr-heights} gives
$h_{K_0}(v_i/\bar v_i)=\tfrac12\log\operatorname N v_i$.
Multiplication by roots of unity and inversion preserve height.
This proves \eqref{dc-exact-height}, including unit residuals.

For the branch calculation,
$G(X)/\zeta=X^2+(2-\zeta)X+1$ has discriminant $-1-3\zeta$
of norm thirteen. Thus $G$ has two distinct roots, as does $\bar G$.
They have no common root: their difference is
$(\zeta-\bar\zeta)N(X)$, whereas at either root of $N$ their
value is $X\ne0$. Their four roots are also disjoint from
$\beta_+,\beta_-$. At infinity $G/\bar G$ tends to the nonzero
constant $\zeta/\bar\zeta$.

Over $\overline{K_0}$ the second equation therefore defines a
cyclic cover of degree $g$, with exactly four branch points,
each of inertia $g$. Simple zeros and poles give this conclusion
also for composite $g$. The first equation gives a cyclic cover
of degree $h$, with exactly the two branch points
$\beta_+,\beta_-$, each of inertia $h$. Both are unramified
at infinity.

The two geometric Galois extensions have disjoint branch sets.
Their intersection is everywhere unramified over $\mathbb P^1$,
and is trivial by Riemann--Hurwitz. Their compositum is thus a
field of degree $hg$, proving geometric connectedness for arbitrary
$\gcd(h,g)$. Roots of unity are only adjoined for this geometric
calculation, not to the arithmetic field $K_0$. Summing the six
ramification contributions gives
\[
 2G_C-2=-2D+2D(1-1/h)+4D(1-1/g),
\]
and hence \eqref{dc-genus}.

The actual positive $x$ avoids the roots of $N$. Also
$G(x)\bar G(x)=F/b^4>0$, so it avoids the other four branch
values. Its nonzero coordinates give a smooth affine point, since
the derivatives in $Y,Z$ are nonzero in characteristic zero.
It lifts to the normalization over $K_0$. No converse asserting
that every point is a primitive integer seed is used.
\end{proof}

At $h=g=2$ the genus is three, consistently with the earlier
two-square-class geometry. This does not assert a seed in the
already excluded subbranch with two unit residuals and even exponents.

\begin{theorem}[Two cumulative residual counts]\label{dc-count-height}
Fix $h,g$. All actual profiles with $V_0\le X_0$, $V_1\le X_1$,
where $X_0,X_1\ge1$, lift to a list of at most
\begin{equation}\label{dc-cover-count}
 CX_0X_1
\end{equation}
models, for an absolute constant $C$. The models can be written as
\begin{equation}\label{dc-polynomial-model}
 (x-\beta_-)Y^h-\eta_0(x-\beta_+)=0,\qquad
 \bar G(x)Z^g-\eta_1G(x)=0.
\end{equation}
Their maximum individual coefficient height, and their projective
coefficient height up to an absolute factor, are bounded by
\begin{equation}\label{dc-model-height}
 H_{\rm coef}=C_0(1+\log V_0+\log V_1),
\end{equation}
where $C_0$ is fixed. The degrees $h,g$ remain explicit and are
not included in this coefficient-height assertion.
\end{theorem}
\begin{proof}
The lattice count of Lemma~\ref{mx-actual-coefficient} bounds the
number of Eisenstein integers of norm at most $X_i$ by $16X_i$.
There are only six choices for each unit and two choices for $e$.
Multiplying these two cumulative counts proves
\eqref{dc-cover-count}. Restricting to compatible actual
coefficients can only reduce the list. No additional sum over
residual norms or independent unit-class factor is added.

The coefficients of $G,\bar G$ and the branch values are fixed.
The product-height inequality and \eqref{dc-exact-height} give
\eqref{dc-model-height}. There are a fixed number of nonzero
coefficients, independent of $h,g$, so their projective height
is bounded by a constant multiple of the same expression.
\end{proof}

For any $L_0,L_1\ge0$, the union with
$V_0\le\exp(L_0h)$ and $V_1\le\exp(L_1g)$ consequently has size
at most
\begin{equation}\label{dc-small-residual-count}
 C\exp(L_0h+L_1g).
\end{equation}
This uses the actual quadratic-norm structure of $F$. It does not
provide small representatives for every class in the generic QC
list. The difficult unit classes in Lemma~\ref{qc-unit-lower-bound}
were not proved to have actual seeds, so their lower bound is fully
consistent with the present coefficient construction.

\begin{corollary}[Height separation with unrestricted exponent ratios]
\label{dc-height-separation}
Write $H=\max(a,b)$ and $H_{\rm point}=\log H$. Every actual
profile satisfies
\begin{equation}\label{dc-point-lower}
 H_{\rm point}\ge
 \max\left\{\frac h2\log7-\log2,
             \frac g4\log7-\frac14\log13\right\}.
\end{equation}
For any sequence of actual profiles such that
\begin{equation}\label{dc-height-domain}
 T=\max(h,g)\longrightarrow\infty,\qquad
 \log V_0+\log V_1=o(T),
\end{equation}
the selected models obey
\begin{equation}\label{dc-height-limit}
 \frac{H_{\rm coef}+\log h+\log g}{H_{\rm point}}
 \longrightarrow0.
\end{equation}
In particular this holds whenever
\begin{equation}\label{dc-lambda-domain}
 \lambda_0=\frac{\max\{1,\log V_0\}}h\longrightarrow0,
 \qquad
 \lambda_1=\frac{\max\{1,\log V_1\}}g\longrightarrow0,
\end{equation}
without a condition on $h/g$ or on the common divisors of $h,g$.
For the combined list in \eqref{dc-small-residual-count} with
$L_i=\lambda_i$, its logarithmic size divided by $T$ also tends to zero.
\end{corollary}
\begin{proof}
We have $M\ge7^h$, $M<c^2$ and $H\ge c/2$.
Also $F\ge7^g$ and $F\le13H^4$. These imply
\eqref{dc-point-lower}, whose right side is at least
$\tfrac14T\log7-\log2$. Under \eqref{dc-height-domain},
\eqref{dc-model-height} is $o(T)$ and
$\log h+\log g\le2\log T=o(T)$. Division proves
\eqref{dc-height-limit}.

Under \eqref{dc-lambda-domain}, both exponents tend to infinity and
\[
 \frac{\log V_0+\log V_1}{T}\le\lambda_0+\lambda_1\longrightarrow0.
\]
Finally the logarithm of \eqref{dc-small-residual-count}, divided
by $T$, is at most $\log C/T+\lambda_0+\lambda_1$.
\end{proof}

These conclusions are necessary properties of an actual sequence,
not a construction of one. High points on varying high-genus curves
are compatible with all the displayed estimates. A uniform bound
\[
 H_{\rm point}\le A(H_{\rm coef}+\log h+\log g+1)
\]
for actual lifted points, with $A$ independent of the profiles, would
exclude \eqref{dc-height-domain}. No such upper bound is proved here;
coefficient height, genus and the number of twists alone do not give it.
Moving roots, nonunit residuals, and the fixed positive two-step
threshold remain unresolved. The present result is an independently
reviewed ordinary proof; its number-field, geometric and height
arguments are not claimed as Lean formalizations.
